%! Author = sriadityavedantam
%! Date = 3/25/26

\section{Consequences of Completeness}

\lesson{4}{Wednesday 20 August 2025 9:10}{Day 4}

\begin{theorem}[Complete Exponential Theorem]{
    Define $S\coloneqq\{x\in\r : x^n < c\}$, $S$ is bounded above, then set $\alpha \coloneqq \sup(S)$ and $\alpha^n = c$.
}
    Suppose $\alpha^n < c$.
    Then there exists $\epsilon > 0$ such that $(\alpha + \epsilon)^n < c$.
    Then $\alpha + \epsilon \in S$, contradicting that $\alpha$ is an upper bound.
    Suppose $\alpha^n > c$.
    Then there exists $\epsilon > 0$ such that $(\alpha - \epsilon)^n > c$.
    Then $\alpha - \epsilon$ is still an upper bound of $S$ smaller than $\alpha$, contradicting minimality.
    Hence $\alpha^n = c$.
\end{theorem}

\begin{definition}[Cardinality of Sets]
    It is the number of elements in set $S$.
\end{definition}

\begin{definition}[Equal Cardinality]
    Sets have the same cardinality if there is a bijective function $f: A\leftrightarrow B$.
\end{definition}

\begin{definition}[Countably Infinite]
    A set is countably infinite if it has the same cardinality as $\n$.
\end{definition}

\begin{definition}[Countable]
    A set is countable if it is either finite, or countably infinite.
    Set $S$ is also countable if its elements can be listed as a sequence.
\end{definition}

\begin{theorem}{
    $\z,\q$ are both countably infinite.
}
    $\z$ -- Define
    \begin{align*}
        f(2n) &= n\\
        f(2n-1) &= -n.
    \end{align*}
    This defines a bijection $f:\n\to\z$.

    $\q$ -- Enumerate $\q$ by arranging fractions $\frac{a}{b}$ in a grid and traversing diagonally, skipping duplicates.
\end{theorem}

\begin{theorem}[Nested Interval Theorem]{
    Suppose $[a_1,b_1], [a_2,b_2],\ldots$ is a sequence of closed bounded intervals in $\r$, which is nested.
    \[
        [a_n,b_n]\supset [a_{n+1},b_{n+1}], \forall n\in\n
    \]
    Then:
    \[
        \bigcap\limits_{i=1}^{\infty} [a_i,b_i] \neq \varnothing
    \]
}
    Let $A\coloneqq \{a_1,a_2,\ldots\},\ B\coloneqq \{b_1,b_2,\ldots\}$.
    For any $b\in B$, $b$ is an upper bound of $A$.
    Thus $\sup(A)\leq\inf(B)$.
    Since $\sup(A)\geq a_i$ for all $i$ and $\sup(A)\leq b_i$ for all $i$, we have $\sup(A)\in[a_i,b_i]$ for all $i$.
    Hence $\sup(A)\in \bigcap_{i=1}^\infty [a_i,b_i]$.
\end{theorem}

\begin{remark}
    The direct statement of this corollary is:
    \begin{displayquote}
        If $[a_1,b_1], [a_2,b_2],\ldots$ is a sequence of closed bounded intervals in $\r$ is nested, then $\bigcap_{i=1}^{\infty} [a_i,b_i]\neq\varnothing$.
    \end{displayquote}
    A contrapositive of this lemma is:
    \begin{displayquote}
        If $\bigcap_{i=1}^{\infty} [a_i,b_i] = \varnothing$, then the intervals are not nested.
    \end{displayquote}
\end{remark}

\section{Cardinality}

\begin{theorem}{
    $\r$ is uncountable.
}
    Suppose $\r$ is countable, so $\r=\{x_1,x_2,x_3,\ldots\}$.
    Construct nested intervals $[a_n,b_n]$ such that $x_n\notin[a_n,b_n]$ for each $n$, and $[a_{n+1},b_{n+1}]\subset[a_n,b_n]$.
    By the Nested Interval Theorem,
    \[
        \bigcap_{n=1}^\infty [a_n,b_n]\neq \varnothing.
    \]
    Let $x$ be in the intersection.
    Then $x\neq x_n$ for all $n$, contradicting that the sequence lists all real numbers.
    Hence, $\r$ is uncountable.
\end{theorem}

\begin{theorem*}{
    If $A\subseteq B$ and $B$ is countable, then $A$ is either countable or finite.
}
\end{theorem*}

\begin{remark}
    A contrapositive of this lemma is:
    \begin{displayquote}
        If $A$ is uncountable, then $A\nsubseteq B$ or $B$ is uncountable.
    \end{displayquote}
\end{remark}

\begin{theorem*}{
    If $A_i$ are each countable sets for $i\in\n$, then $\bigcup_{i\in\n} A_i$ is countable.\label{thm:1.5.8}
}
\end{theorem*}

\begin{remark}
    A contrapositive of this lemma is:
    \begin{displayquote}
        If $\bigcup_{i\in\n} A_i$ is uncountable, there exists an $A_i$ that is uncountable.
    \end{displayquote}
\end{remark}
